\chapter{Introduction}

\section{Motivation}Identifying highly influential structures within large networks is a fundamental problem in graph analytics.
The \textbf{Betweenness Centrality Problem (BCP)} formalizes this challenge by calculating the betweenness centrality for the nodes of a graph.
This metric quantifies the importance of a node based on the number of shortest paths that pass through it.

Solving the BCP is critical for a wide range of real-world applications.
These include identifying critical vulnerabilities in power grids, optimizing data traffic in telecommunication networks,
mitigating urban traffic congestion, and securing global supply chains against systemic bottlenecks.

Exact algorithms for the BCP traditionally rely on modern shortest-path frameworks,
most notably Brandes' algorithm \parencite{brandes_faster_2001}, which executes multiple graph traversals (such as Dijkstra's algorithm).
Although these exact methods guarantee optimal solutions, their computational runtime becomes prohibitive on large-scale graphs.
Consequently, approximation methods such as \textbf{KADABRA} have been developed, successfully trading strict optimality for significantly
improved efficiency \parencite{borassi_kadabra_2019}.

Recent advances in Graph Neural Networks (GNNs) show great promise.
Specifically, the creators of BRAVA-GNN \parencite{noauthor_degree-mass_2026} have achieved state-of-the-art performance among GNN approaches,
even outperforming KADABRA in full-graph evaluations.
However, existing benchmarks primarily compare total graph calculations.
A primary advantage of KADABRA is its ability to approximate betweenness centrality exclusively for the top-$k$ nodes,
offering substantial speedups by bypassing the rest of the network.
This report extends the comparative analysis between KADABRA and BRAVA-GNN by focusing specifically on this top-$k$ approximation
performance. Furthermore, this report serves as supplementary documentation for the implementation and integration of the KADABRA
algorithm into the Julia \texttt{Graphs.jl} ecosystem, as well as comparing the authors C++ implementation vs this new Julia.

\section[Related Work]{Related Work}\label{sec:related-work}

The development of polynomial-time exact algorithms for DSP has been a significant area of research.
\textcite{picard_network_1982} published the first exact algorithm, which was fundamentally built on maximum-flow computations.
This was later advanced by \textcite{goldberg_finding_1984}, who proposed a version with improved time complexity.
A different approach was taken by \textcite{charikar_greedy_2000}, who designed an exact algorithm based on a linear program (LP).
The existence of an LP-based algorithm is not surprising, as the max-flow problem can be formulated as an LP.
However, Charikar's method offers a new perspective, allowing his greedy-peeling approximation algorithm to be understood as a primal-dual algorithm for DSP.
This greedy algorithm, which iteratively removes nodes with the lowest degree, provides a 2-approximation.

Building upon this concept, \textbf{iterative peeling algorithms} have been developed.
Instead of performing the peeling process once using node degree as the criterion, it is executed multiple times.
In each new iteration, more complex node weights, influenced by the results of the previous iterations, are used instead of the simple degree.
\textcite{boob_flowless_2020} introduced the first algorithm of this type, called \textit{Greedy++},
and demonstrated that it achieves a near-optimal result in just a few iterations and is an order of magnitude faster than Goldberg's maximum flow algorithm.
\textcite{chekuri_densest_2022} showed that Greedy++ is an approximation scheme that can achieve a value arbitrarily close to the optimum given enough iterations,
although the required number of iterations can be extremely high in theory.

\textcite{hochbaum_flow_2024} introduced the \textit{Incremental Parametric Cut} procedure (IPC), which uses the pseudoflow algorithm\textcite{hochbaum_pseudoflow_2008}.
This approach uses previous results for initialization when the density parameter changes monotonically, thereby avoiding wasteful recomputations.
Techniques like \textbf{restartable maximum flow} allow an existing flow solution to be efficiently updated after minor changes to the graph, instead of restarting the algorithm from scratch.
\textcite{hochbaum_flow_2024} claim that their algorithm is faster by several orders of magnitude than \textit{Greedy++}.
However, the publication cited to support this claim is not available online.

Recent advances have produced new parametric flow algorithms with unique capabilities.
The Parametric Breadth-First Search (PBFS) \parencite{beines_simpler_2024} and Parametric Pseudo Flow (PPF) \parencite{Sauer_PPF_2025}
are the first algorithms designed to find a problem's breakpoints—critical parameter values where the minimum cut's structure changes—in strictly monotone order.
This specific property opens up new strategies for solving ratio problems like the DSP.

The optimal density for the DSP corresponds to the breakpoint of the highest value.
The existence of algorithms that find breakpoints in order enables a new solution strategy:
if the problem can be reformulated such that this highest-value breakpoint becomes the \textit{first} to be found, the solution process could be significantly accelerated.

\section[Our Contribution]{Our Contribution}\label{sec:our-contribution}


A central contribution of this work is the first systematic, empirical comparative study of modern parametric maximum-flow algorithms for the densest subgraph problem (DSP)
versus established approximation methods. The investigation covers 47 real-world graphs and includes the first comprehensive and publicly documented evaluation of the
Incremental Parametric Cut procedure (IPC) by \textcite{hochbaum_flow_2024}, given that their own empirical study is not publicly available.
%This thesis makes several contributions to the understanding and practical application of algorithms for the Densest Subgraph Problem.
%A central contribution of this thesis is the first systematic, empirical comparative study of modern parametric maximum flow algorithms for DSP against established approximation methods.
%, which spans 47 real-world graphs and includes the first comprehensive and publicly documented analysis of the
%\textit{Incremental Parametric Cut} (IPC) by \textcite{hochbaum_flow_2024}, considering that their own empirical study is not publicly available.
To enable this evaluation we implement and assess two novel exact algorithms:

\textbf{{\DSPwithFBA}}, which leverages recent monotone parametric flow algorithms to find the solution directly by identifying the first breakpoint in a reformulated network, making it the only one required.

\textbf{{\DSPwithR}}, an intersection-based algorithm using restartable maximum flow with both push-relabel and pseudoflow implementations, a similar approach to IPC.

Our experiments reveal that the choice of the best algorithm strongly depends on graph properties, particularly degree variance.
For graphs with low variance, \textbf{DSP-FBA} performs best, whereas for graphs with high variance, IPC is superior.
Confirming the hypothesis previously stated by Hochbaum (2024), our results challenge the common assumption that approximation algorithms are always preferable for large graphs.
We demonstrate that exact algorithms can achieve optimal solutions with runtimes only slightly higher than those of approximations like Greedy++,
making them a viable alternative when solution quality is critical.

Additionally, we systematically analyze the effectiveness of graph pruning based on the greedy algorithm by Charikar (2000).
We show that pruning offers substantial benefits, especially for unweighted graphs with high degree variance, and support this with more extensive experiments.
Specifically, we investigate the combination of pruning followed by these fast flow algorithms—a configuration not previously analyzed.
We show that with pruning, exact algorithms on unweighted graphs can consistently outperform the unpruned Greedy++ approximation, while the benefit is far less pronounced for weighted graphs.



Our implementation extends the parametric maximum flow framework by \textcite{beines_simpler_2024} with new, DSP-specific algorithms.
These contributions provide both theoretical insights into the performance of exact versus approximate algorithms and practical guidance for practitioners.



